% =====================================================================
%  Questions / réponses — version vérifiée
% =====================================================================
\documentclass[a4paper,11pt]{article}
% =====================================================================
%  Préambule commun aux documents de soutenance
% =====================================================================
\usepackage[T1]{fontenc}
\usepackage[utf8]{inputenc}
\usepackage[sfdefault]{carlito}
\usepackage[french]{babel}
\usepackage{csquotes}
\usepackage[a4paper,margin=2.2cm,top=2.5cm,bottom=2.2cm]{geometry}
\usepackage{microtype}
\usepackage[htt]{hyphenat}
\tolerance=2000
\emergencystretch=3em
\widowpenalty=10000
\clubpenalty=10000
\setlength{\parskip}{5pt plus 2pt minus 1pt}

\usepackage{xcolor}
\usepackage{colortbl}
\usepackage{array}
\usepackage{tabularx}
\usepackage{longtable}
\usepackage{enumitem}
\usepackage{ragged2e}
\usepackage{titlesec}
\usepackage{fancyhdr}
\usepackage{tcolorbox}
\usepackage[hidelinks]{hyperref}

\definecolor{btkblue}{HTML}{14507A}
\definecolor{btklight}{HTML}{1B6CA8}
\definecolor{btkpale}{HTML}{EAF3FA}
\definecolor{btkambre}{HTML}{9A6B12}
\definecolor{btkambrepale}{HTML}{FDF6E6}
\definecolor{btkvert}{HTML}{2E7D5B}
\definecolor{btkvertpale}{HTML}{EDF7F2}

\renewcommand{\tabularxcolumn}[1]{>{\RaggedRight\arraybackslash}p{#1}}
\newcolumntype{L}[1]{>{\RaggedRight\arraybackslash}p{#1}}
\renewcommand{\arraystretch}{1.25}

\titleformat{\section}{\Large\bfseries\color{btkblue}}{\thesection}{0.7em}{}
\titleformat{\subsection}{\large\bfseries\color{btklight}}{\thesubsection}{0.6em}{}
\titlespacing*{\section}{0pt}{2.4ex plus 1ex}{1.2ex}
\titlespacing*{\subsection}{0pt}{1.8ex plus .8ex}{0.8ex}

\pagestyle{fancy}
\fancyhf{}
\renewcommand{\headrulewidth}{0.4pt}
\renewcommand{\headrule}{\hbox to\headwidth{\color{btklight}\leaders\hrule height \headrulewidth\hfill}}
\fancyhead[R]{\small\color{btkblue}Oussema Korbosli --- PFE BTK}
\fancyfoot[C]{\small\thepage}

% Encadré « à savoir par cœur »
\newtcolorbox{coeur}[1][]{
  colback=btkpale, colframe=btklight, boxrule=0.9pt, arc=2pt,
  left=7pt, right=7pt, top=5pt, bottom=5pt,
  fonttitle=\bfseries\small, title={À savoir par c\oe{}ur}, #1}

% Encadré « note » : conseil de présentation
\newtcolorbox{note}[1][]{
  colback=btkvertpale, colframe=btkvert, boxrule=0.8pt, arc=2pt,
  left=7pt, right=7pt, top=4pt, bottom=4pt, #1}

% Encadré « chiffre corrigé »
\newtcolorbox{corrige}[1][]{
  colback=btkambrepale, colframe=btkambre, boxrule=0.9pt, arc=2pt,
  left=7pt, right=7pt, top=5pt, bottom=5pt,
  fonttitle=\bfseries\small, title={Chiffre corrigé}, #1}

% Question / réponse
\newcommand{\question}[1]{\subsection{#1}}
\newcommand{\cle}[1]{\par\vspace{2pt}\noindent
  \colorbox{btkpale}{\parbox{\dimexpr\linewidth-2\fboxsep}{%
  \textbf{Phrase à dire :} \emph{#1}}}\par\vspace{3pt}}

\fancyhead[L]{\small\color{btkblue}Questions / réponses}

\begin{document}

\begin{center}
{\color{btkblue}\rule{\textwidth}{2pt}}\\[0.3cm]
{\huge\bfseries\color{btkblue} Questions et réponses\par}
\vspace{0.2cm}
{\large Préparation de la soutenance --- PFE BTK\par}
\vspace{0.2cm}
{\color{btkblue}\rule{\textwidth}{2pt}}
\end{center}

\vspace{0.4cm}
\begin{corrige}[title={Version vérifiée --- ce qui a changé}]
Quatre corrections de chiffres (Q5, Q7, Q8), une question en double supprimée
(Q3/Q4 étaient identiques), et le tableau des tests remplacé par celui du
rapport --- l'ancien citait une suppression de client que l'application ne sait
pas faire. Récapitulatif en dernière page.
\end{corrige}

% =====================================================================
\section{Architecture}

\question{Q1 --- Est-ce que c'est du MVC ?}

\begin{itemize}[leftmargin=*,itemsep=1pt]
  \item Non : c'est une architecture \textbf{en couches} (N-tier), pas un MVC pur.
  \item Le MVC classique compte trois composants : Model, View, Controller.
  \item Ici il y a \textbf{quatre couches} : Présentation, Contrôle, Métier, Persistance.
  \item Dans un MVC pur, \enquote{Métier} et \enquote{Persistance} seraient fondus dans le Model.
  \item Le Model a donc été découpé en deux sous-couches, pour affiner la séparation.
\end{itemize}
\cle{C'est une architecture en couches inspirée du MVC : Présentation = View,
Contrôle = Controller, et le Model est affiné en Métier + Persistance.}

\question{Q2 --- Pourquoi pas un MVC pur, alors ?}

\begin{itemize}[leftmargin=*,itemsep=1pt]
  \item \textbf{Responsabilité unique} : le Model du MVC classique mélange règles métier et accès aux données.
  \item \textbf{Testabilité} : la couche Métier peut être testée seule, sans la vraie base Oracle.
  \item \textbf{Cohérence technique} : c'est le standard Jakarta EE --- EJB pour le métier, JPA pour la persistance.
  \item \textbf{Évolutivité} : si la base change, ou si on ajoute une API REST, une seule couche est touchée.
\end{itemize}
\cle{On n'a pas voulu un MVC pur pour éviter de surcharger le Model avec deux
responsabilités différentes. En les séparant, on gagne en clarté, en
testabilité et en flexibilité.}

\question{Q3 --- Différence entre méthode et méthodologie ?}

\begin{itemize}[leftmargin=*,itemsep=1pt]
  \item \textbf{Méthode} : une pratique concrète, un ensemble d'outils et de techniques appliqués --- par exemple Scrum.
  \item \textbf{Méthodologie} : le cadre conceptuel qui guide le choix et l'enchaînement des méthodes.
\end{itemize}
\cle{La méthodologie est le cadre théorique, la méthode est son application
concrète --- Scrum est une méthode issue de la méthodologie agile.}

\begin{corrige}
Cette question figurait \textbf{deux fois} (Q3 et Q4), avec exactement la même
réponse. Le doublon est supprimé, et la place est prise par une question qui
tombe souvent : le choix de CRISP-DM.
\end{corrige}

\question{Q4 --- Pourquoi CRISP-DM et pas Scrum ?}

\begin{itemize}[leftmargin=*,itemsep=1pt]
  \item Le projet est \textbf{centré sur les données}, pas seulement sur la production de code.
  \item Scrum organise la fabrication logicielle mais ne dit rien de la compréhension des données ni de leur valorisation.
  \item CRISP-DM est la méthodologie de référence des projets décisionnels : compréhension métier, compréhension des données, préparation, modélisation, évaluation, déploiement.
  \item Ses six phases épousent le déroulement réel du projet, tout en gardant des pratiques agiles : itérations et validations régulières.
\end{itemize}
\cle{Scrum organise la production logicielle ; CRISP-DM organise le travail sur
la donnée. Mon projet allant de la donnée brute jusqu'à la décision, c'est
CRISP-DM qui décrit le mieux mon déroulement.}

\question{Q9 --- Pourquoi Jakarta EE et pas Spring Boot, plus répandu ?}

\begin{itemize}[leftmargin=*,itemsep=1pt]
  \item \textbf{Standard robuste et éprouvé}, avec des spécifications stables : Servlets, EJB, JPA, CDI.
  \item \textbf{Portabilité} : aucune dépendance à un serveur précis.
  \item \textbf{Cohérence avec l'environnement bancaire}, où les serveurs d'applications de ce type sont déjà en place.
  \item \textbf{Séparation des responsabilités native} : EJB pour le métier, JPA pour la persistance --- ce qui s'aligne exactement sur l'architecture en couches retenue.
\end{itemize}
\cle{J'ai choisi Jakarta EE parce que c'est un standard robuste et portable,
cohérent avec un environnement bancaire où la stabilité et la conformité aux
standards priment sur la rapidité de développement.}

\question{Q10 --- Que voulez-vous dire par \enquote{workflows de validation} ?}

\begin{itemize}[leftmargin=*,itemsep=1pt]
  \item Certaines actions sensibles ne s'appliquent pas directement : elles passent par un circuit d'approbation.
  \item Exemple : le directeur commercial soumet une demande de création d'agence, l'administrateur la valide avant qu'elle prenne effet.
  \item Une modification de données personnelles va plus loin : \textbf{double validation} --- le directeur approuve, puis l'administrateur applique.
  \item Cela évite qu'une seule personne modifie une donnée critique sans second contrôle.
\end{itemize}
\cle{Un workflow de validation, c'est un circuit d'approbation : une action
sensible est d'abord soumise, puis validée par un autre rôle avant d'être
appliquée. Aucune donnée critique n'est modifiée par une seule personne.}

% =====================================================================
\section{Segmentation et intelligence artificielle}

\question{Q5 --- Pourquoi K-Means et pas un autre algorithme ?}

\begin{itemize}[leftmargin=*,itemsep=1pt]
  \item J'ai comparé quatre modèles : K-Means, agglomératif, GMM et DBSCAN.
  \item K-Means obtient le \textbf{meilleur score de silhouette : 0,604}, et le meilleur Calinski-Harabasz : 95,2.
  \item DBSCAN a un Davies-Bouldin plus favorable (0,393 contre 0,709), \textbf{mais il refuse de classer CENTRALE}, qu'il traite comme du bruit.
  \item Or un modèle de pilotage doit affecter \textbf{toutes} les agences à un segment : je retiens donc K-Means.
  \item Il est en outre simple, rapide, et permet de classer facilement une nouvelle agence.
\end{itemize}
\cle{J'ai comparé quatre algorithmes. K-Means a la meilleure silhouette.
DBSCAN a un meilleur Davies-Bouldin mais il rejette une agence comme bruit :
je retiens le modèle qui classe toutes les agences, condition pour s'en servir
en pilotage.}

\begin{corrige}
La version d'origine annonçait \enquote{silhouette 0,736 et Davies-Bouldin
0,350, les meilleurs des quatre}. Ces valeurs venaient de l'ancienne exécution
sur un jeu synthétique. Sur les données réelles : \textbf{0,604} et
\textbf{0,709}, et le meilleur Davies-Bouldin est celui de DBSCAN. Le tableau
du rapport porte les valeurs réelles.
\end{corrige}

\noindent\begin{tabularx}{\textwidth}{|L{3.4cm}|c|c|c|X|}
\hline
\rowcolor{btkpale}
\textbf{Modèle} & \textbf{k} & \textbf{Silhouette} & \textbf{D.-Bouldin} & \textbf{Remarque}\\
\hline
K-Means & 3 & \textbf{0,604} & 0,709 & retenu\\
\hline
DBSCAN & 4 & 0,517 & \textbf{0,393} & rejette CENTRALE comme bruit\\
\hline
GMM & 3 & 0,446 & 0,770 & \\
\hline
Agglomératif & 3 & 0,417 & 0,841 & \\
\hline
\end{tabularx}

\question{Q6 --- Pourquoi k = 3 ?}

\begin{itemize}[leftmargin=*,itemsep=1pt]
  \item Deux méthodes combinées : la \textbf{méthode du coude} et le \textbf{score de silhouette}.
  \item La silhouette est en réalité maximale à k = 2 --- mais k = 2 ne fait qu'isoler trivialement les entités sans activité.
  \item J'ai donc retenu la meilleure solution \textbf{à partir de k = 3}, la première qui soit exploitable pour le pilotage.
\end{itemize}
\cle{k = 2 sépare trivialement les entités sans activité. J'ai retenu la
meilleure solution à partir de k = 3, qui donne trois profils interprétables et
utilisables par la direction.}

\begin{corrige}
La version d'origine disait que \enquote{les deux méthodes convergent vers 3}.
C'est inexact et vérifiable dans la courbe de silhouette du rapport : le maximum
est à k = 2 (0,681), k = 3 vient ensuite (0,604). Mieux vaut l'annoncer
soi-même avec la justification métier.
\end{corrige}

\question{Q7 --- Pourquoi standardiser les données avant K-Means ?}

\begin{itemize}[leftmargin=*,itemsep=1pt]
  \item K-Means repose sur la \textbf{distance euclidienne}.
  \item Sans standardisation, une variable à grande échelle --- la production de crédits, en centaines de millions --- écraserait toutes les autres.
  \item \texttt{StandardScaler} centre et réduit : moyenne 0, écart-type 1, toutes les variables au même poids.
  \item J'ai en plus appliqué une \textbf{transformation logarithmique} aux montants et au portefeuille, dont la distribution est très asymétrique.
\end{itemize}
\cle{K-Means utilise la distance euclidienne : sans standardiser les six
variables, celles à grande échelle fausseraient les distances. J'ai aussi passé
les montants au logarithme, parce que leur distribution est très asymétrique.}

\begin{corrige}
La version d'origine parlait des \enquote{7 variables}. La segmentation en
utilise \textbf{six} : effectif, gestionnaires, clients, comptes, crédits,
épargne. Le taux de présence n'est pas retenu --- il mesure l'assiduité, pas la
taille ni l'activité commerciale.
\end{corrige}

\question{Q8 --- C'est quoi le score de silhouette ?}

\begin{itemize}[leftmargin=*,itemsep=1pt]
  \item Une mesure de qualité du regroupement, comprise entre $-1$ et $1$.
  \item Proche de 1 : les points sont serrés dans leur groupe et bien séparés des autres.
  \item Ici : \textbf{0,604}, ce qui indique des groupes nettement distincts.
\end{itemize}
\cle{Le score de silhouette évalue si un point est proche de son propre groupe
et loin des autres. Notre 0,604 indique des groupes bien séparés.}

\question{Q16 --- Comment avez-vous réalisé la segmentation ?}

J'ai utilisé le datamart produit par la chaîne ETL : \textbf{49 entités} du
réseau et \textbf{six indicateurs}. Le siège est écarté de l'analyse --- 540
employés, c'est une structure centrale, pas une agence --- ce qui laisse
\textbf{48 entités}. J'ai appliqué une transformation logarithmique aux
montants, standardisé les variables, puis retenu trois groupes avec la méthode
du coude et la silhouette.

\question{Q17 --- Pourquoi avoir standardisé les données ?}

Parce que les variables n'ont pas les mêmes échelles. La standardisation les
rend comparables et évite qu'une variable à grandes valeurs domine les autres
dans le calcul des distances.

\question{Q18 --- Pourquoi avoir choisi K-Means ?}

J'ai comparé quatre modèles. K-Means obtient le meilleur score de silhouette, et
il est le plus pertinent pour le pilotage : chaque agence est affectée à un
segment, et une nouvelle agence peut être classée immédiatement.

\question{Q19 --- Quels modèles avez-vous comparés ?}

K-Means, l'agglomératif, GMM et DBSCAN.

\begin{note}
\small Citer les quatre suffit, à condition de savoir justifier le choix de
K-Means (voir Q5).
\end{note}

% =====================================================================
\section{Décisionnel, KPI et Power BI}

\question{Q11 --- Pourquoi ces KPI, et sous quels critères ?}

Les KPI ont été choisis en fonction des besoins métier identifiés et des données
réellement disponibles dans les sources de la banque. L'objectif était de couvrir
les principaux axes de pilotage du réseau : les ressources humaines, l'activité
commerciale, la clientèle et l'assiduité. J'ai donc retenu des indicateurs qui
permettent à la direction de mesurer la taille et l'activité des agences, leur
performance commerciale et la présence des employés.

\question{Q15 --- Quels sont vos KPI ?}

Le nombre de clients, la production de crédits, le total des comptes ouverts, la
collecte d'épargne, le taux de présence, la répartition des rôles et la
segmentation des agences. Ils évaluent à la fois la performance commerciale des
agences, leur capacité à développer leur portefeuille, et leur organisation
interne.

\begin{note}
\small Si on demande un seul exemple : \emph{nombre de clients} mesure la taille
du portefeuille de l'agence ; \emph{taux de présence} mesure l'assiduité des
équipes.
\end{note}

\question{Q12 --- Comment fonctionne votre chaîne ETL ?}

Elle est implémentée en Python avec pandas, et exécutée dans un notebook
connecté à la base Oracle. Elle enchaîne cinq étapes : collecte des cinq tables
sources, nettoyage, transformation, intégration sur la clé \texttt{SK\_AGENCE},
et contrôle de qualité. Le contrôle final vérifie l'absence de valeur manquante,
de valeur négative et d'agence en double --- et \textbf{interrompt le
chargement} si l'un d'eux échoue.

\question{Q13 --- C'est quoi une mesure DAX ?}

C'est un calcul défini dans Power BI qui produit dynamiquement un indicateur à
partir des données du modèle. Il est évalué \textbf{au moment de l'affichage},
dans le contexte des filtres actifs.

\question{Q14 --- Donnez-moi un exemple de mesure.}

\begin{itemize}[leftmargin=*,itemsep=1pt]
  \item \textbf{Nb Employés} : compte le nombre distinct d'employés à partir de \texttt{SK\_UTILISATEUR}.
  \item \textbf{Nb Clients} : compte le nombre distinct de clients à partir de \texttt{SK\_CLIENT}.
  \item \textbf{Taux Présence} : rapport entre les pointages \enquote{présent} ou \enquote{retard} et le total des pointages.
\end{itemize}

\begin{coeur}
\textbf{Mesure ou colonne calculée ?} C'est la relance la plus fréquente après
une question DAX. Une \textbf{colonne calculée} est évaluée à l'actualisation et
stockée : elle sert à \emph{filtrer ou grouper}. Une \textbf{mesure} est évaluée
à l'affichage, dans le contexte des filtres : elle sert à \emph{agréger}. Dans
mon rapport, \texttt{Rôle} est une colonne calculée, \texttt{Taux Présence} une
mesure.
\end{coeur}

\begin{note}
\small Les données du volet IA proviennent du \textbf{datamart produit par la
chaîne ETL}, et non directement de la base opérationnelle.
\end{note}

% =====================================================================
\section{Les tests réalisés}

\begin{corrige}
Le tableau d'origine listait dix scénarios dont deux ne correspondent pas à
l'application : \enquote{Suppression d'un client} --- l'application ne sait pas
supprimer un client, il n'existe aucune opération de suppression dans le service
correspondant --- et \enquote{Modification d'un employé}, qui ne se fait pas
directement mais passe par le circuit de demande à double validation. Le tableau
ci-dessous est celui du \textbf{rapport remis au jury} : c'est celui qu'il faut
citer, pour ne pas se contredire.
\end{corrige}

\noindent\begin{tabularx}{\textwidth}{|c|>{\bfseries}L{2.6cm}|X|L{2.3cm}|}
\hline
\rowcolor{btkpale}
\textbf{\#} & \textmd{\textbf{Module}} & \textbf{Scénario testé} & \textbf{Résultat}\\
\hline
1 & Authentification & Connexion avec des identifiants valides puis invalides. & Conforme\\
\hline
2 & Rôles & Accès à une fonction réservée à l'\texttt{ADMIN} par un \texttt{USER}. & Conforme\\
\hline
3 & Agences & Création, modification et suppression d'une agence. & Conforme\\
\hline
4 & Employés & Ajout d'un employé et rattachement à une agence. & Conforme\\
\hline
5 & Clients & Affectation d'un client à un gestionnaire. & Conforme\\
\hline
6 & Objectifs & Consultation de la production par agence et période. & Conforme\\
\hline
7 & Pointage & Pointage au-delà de 9~h (détermination du statut). & Conforme (\texttt{RETARD})\\
\hline
8 & Demandes & Validation puis refus d'une demande, avec journalisation. & Conforme\\
\hline
9 & ETL & Exécution du pipeline et production du datamart. & Conforme\\
\hline
10 & Décisionnel & Filtrage par district dans le rapport Power BI. & Conforme\\
\hline
\end{tabularx}

\begin{note}
\small Si on vous demande quels tests vous avez faits, en citer deux ou trois et
décrire le résultat obtenu. Les plus faciles à raconter : le pointage après 9~h
qui bascule en \texttt{RETARD}, et l'accès refusé à un \texttt{USER} sur une
page réservée à l'\texttt{ADMIN}.
\end{note}

% =====================================================================
\newpage
\section*{Récapitulatif des corrections}

\noindent\begin{tabularx}{\textwidth}{|>{\columncolor{btkpale}\bfseries}L{1.6cm}|L{5.4cm}|X|}
\hline
\textmd{\textbf{Point}} & \textbf{Version d'origine} & \textbf{Version vérifiée}\\
\hline
Q3 / Q4 & question posée deux fois, réponse identique & doublon supprimé, remplacé par le choix de CRISP-DM\\
\hline
Q5 & silhouette \textbf{0,736}, Davies-Bouldin \textbf{0,350}, \enquote{les meilleurs} & silhouette \textbf{0,604}, Davies-Bouldin \textbf{0,709} ; DBSCAN a un meilleur Davies-Bouldin mais rejette une agence\\
\hline
Q6 & \enquote{les deux méthodes convergent vers 3} & la silhouette est maximale à k = 2 ; k = 3 est le meilleur choix exploitable\\
\hline
Q7 & \enquote{les \textbf{7} variables} & \textbf{six} variables ; le taux de présence n'est pas retenu\\
\hline
Q8 & \enquote{0,736 = bon résultat} & \textbf{0,604}\\
\hline
Tests & \enquote{Suppression d'un client}, \enquote{Modification d'un employé} & remplacés par les dix scénarios du rapport\\
\hline
\end{tabularx}

\vspace{0.8em}
\noindent\textbf{Inchangé et vérifié :} Q1, Q2, Q9, Q10, Q11, Q12, Q13, Q14,
Q15, Q16 (49 entités et 6 indicateurs : déjà juste), Q17, Q18, Q19.

\end{document}
